\chapter{Обзор литературы}

\label{chap:background}

\section{Протокол языкового сервера}

Протокол языкового сервера (LSP) \cite{language-server-protocol:2016} — это
языко-независимый стандарт, который упрощает интеграцию функциональности языка
программирования в различные среды разработки. Он решает проблему необходимости
создавать отдельную реализацию для каждой пары «язык-редактор».

\subsection*{Реализации языковых серверов}

\label{sec:background:implementations}

Для оптимальной интеграции с компилятором Jasmin, написанным на OCaml, языковой
сервер также был разработан на OCaml. Ниже анализируются релевантные реализации
на OCaml для выявления лучших практик.

\subsubsection*{Языковой сервер Jasmin}

Помимо нашей реализации, существует только один сервер для Jasmin
\cite{gh-jasmin-lsp}, который предоставляет лишь базовую диагностику и обладает
рядом недостатков.

\subsubsection*{Языковой сервер OCaml (Merlin)}

Merlin \cite{gh-merlin} — это продвинутый сервер для OCaml, разработанный до
появления LSP \cite{merlin:2018}. Его ключевое преимущество — качественная
обработка неполного или некорректного кода, что позволяет предоставлять точную
диагностику и автодополнение в реальном времени. Архитектура Merlin является
эталоном глубокой интеграции с IDE.

\subsubsection*{Языковой сервер Semgrep}

Языковой сервер Semgrep \cite{gh-semgrep-lsp} послужил основным техническим
ориентиром благодаря чистому и модульному дизайну. Он использует стандартные
библиотеки \texttt{ocaml-lsp} \cite{opam-ocaml-lsp} для работы с протоколом и
\texttt{lwt} \cite{opam-lwt} для асинхронных операций. Применение чистых функций
упрощает отладку и добавление новой функциональности.

\subsection*{Библиотеки для разработки языковых серверов}

Для разработки была использована библиотека \texttt{ocaml-lsp}
\cite{opam-ocaml-lsp}, широко применяемая в сообществе OCaml. Она обрабатывает
сообщения JSON-RPC 2.0 \cite{jsonrpc:2010}, преобразуя их в типизированные
структуры OCaml с помощью библиотеки \texttt{yojson} \cite{gh-yojson}. Это
позволяет разработчикам сосредоточиться на логике языка, а не на деталях
протокола.

\section{Компилятор Jasmin}

Компилятор Jasmin \cite{jasmin:2017} — это формально верифицированный инструмент
для компиляции криптографического кода в безопасный ассемблерный код. Встроенный
статический анализатор проверяет такие свойства, как безопасность памяти и
завершаемость. Компилятор может экспортировать программы в спецификации
EasyCrypt \cite{easycrypt:2013} для дальнейшей формальной верификации. Основная
сложность при работе с ним — отсутствие публичного API.

\section{Tree-sitter}

Tree-sitter \cite{tree-sitter} — это генератор парсеров для построения быстрых и
отказоустойчивых синтаксических анализаторов, которые создают конкретные
синтаксические деревья (CST). Его ключевое преимущество — инкрементальный
парсинг, который позволяет эффективно обновлять синтаксическое дерево при
небольших изменениях в коде, что крайне важно для текстовых редакторов и IDE.